\section{Tiên đề cặp}

\ 

Tiên đề này cho phép chúng ta tạo ra một tập hợp mới từ hai tập hợp đã cho. Nó phát biểu rằng với mọi tập hợp $x$ và $y$, tồn tại một tập hợp $z$ sao cho $z$ chỉ chứa $x$ và $y$ làm phần tử của nó. Phát biểu dưới dạng ngôn ngữ lô-gích:
\begin{equation*}
    \defMath{\forall x, \forall y, \exists z, \forall w \left( w \in z \iff (w = x \lor w = y) \right)}.
\end{equation*}

Có thể thấy chỉ tồn tại một $z$ duy nhất theo tiên đề ngoại diên như khi chúng ta đã chứng minh về tính duy nhất của $\varnothing$. Giả sử tồn tại $z_1$ và $z_2$ thỏa mãn tiên đề cặp. Theo chính tiên đề này:
\begin{equation*}
    \begin{cases}
        a \in z_1 \iff (a = x \lor a = y) \\
        a \in z_2 \iff (a = x \lor a = y)
    \end{cases}.
\end{equation*}
Theo tính chất bắc cầu của phép đẳng giá, chúng ta có:
$$
a \in z_1 \iff a \in z_2
$$
đúng với mọi $a$. Theo tiên đề ngoại diên, $z_1 = z_2$.

Kí hiệu tập này là $\defMath{\{x; y\}}$. Với một lập luận tương đương, dễ dàng thấy rằng $\defMath{\{x; y\} = \{y; x\}}$. Khi này, chúng ta có một \defText{đôi không thứ tự} hay \defText{cặp không thứ tự}. Một điểm nữa cần phải để ý, tiên đề không yêu cầu $x \neq y$, nên chúng ta hoàn toàn có thể xây dựng tập $\{x; x\}$ mà chỉ có một phần tử $x$. Đây là \defText{tập hợp đơn}, và có thể viết tắt tập này là $\defMath{\{x\}}$.

Từ một đôi không thứ tự, chúng ta có thể xây dựng được \defText{đôi có thứ tự} (hoặc \defText{cặp có thứ tự}) theo định nghĩa của Ku-ra-tốp-xki\footnote{Kazimierz Kuratowski (1896 -- 1980).} vào năm 1921: $$\forall x, \forall y,\defMath{(x; y) = \{\{x; y\}; \{x\}\}}.$$
Ku-ra-tốp-xki không phải là người đầu tiên đưa định nghĩa về đôi có thứ tự, nhưng là định nghĩa này trở nên phổ biến do tính đơn giản và hiệu quả của nó. Cụ thể, định nghĩa này đảm bảo rằng:
\begin{quotation}
    Với mọi tập hợp $x$, $y$, $u$, $v$, có $(x; y) = (u; v)$ nếu và chỉ nếu $x = u$ và $y = v$.
\end{quotation}
Để chứng minh điều này, chúng ta chia làm ba trường hợp:

\textcolor{colorEmphasisCyan}{Trường hợp 1 --- $x = y$}: Khi này, $$(x; y) = (x; x) = \{\{x; x\}; \{x\}\} = \{\{x\}; \{x\}\} = \{\{x\}\}.$$
Ngoài ra, còn có $(x; y) = (u; v)$ và định nghĩa $(u; v) = \{\{u; v\}; \{u\}\}$. Điều này dẫn đến 
$$
\{\{x\}\} = \{\{u; v\}; \{u\}\}.
$$
Tiên đề ngoại diên cho $\{x\} = \{u; v\}$. Sử dụng tiên đề ngoại diên một lần nữa để có $x = u = v$. Theo tính chất bắc cầu, chúng ta có $x = u = y = v$.

\textcolor{colorEmphasis}{Trường hợp 2 --- $u = v$}: Lập luận tương tự như trường hợp 1 để có điều tương đương.

\textcolor{colorEmphasisGreen}{Trường hợp 3 --- $x \neq y$ và $u \neq v$}: Do $\{x\} \in (x; y)$ và $(x; y) = (u; v)$ nên $\{x\} \in (u; v)$. Từ đó, $\{x\} = \{u; v\}$ hoặc $\{x\} = \{u\}$ theo tiên đề hợp. Tuy nhiên, $\{x\} = \{u; v\}$ dẫn đến $x = u = v$, mâu thuẫn với giả thiết $u \neq v$. Vậy $\{x\} = \{u\}$, tức là $x = u$ theo tiên đề ngoại diên.

$(x; y) = (u; v)$ cũng cho chúng ta $\{x; y\} \in (u; v)$, hay $\{x; y\} = \{u; v\}$ hoặc $\{x; y\} = \{u\}$. Nhưng nếu $\{x; y\} = \{u\}$, thì, sử dụng tiên đề ngoại diên, $x = y = u$, mâu thuẫn với giả thiết $x \neq y$. Vậy $\{x; y\} = \{u; v\}$, tức là $y = u$ hoặc $y = v$; nhưng nếu $y = u$ thì $x = y = u$, mâu thuẫn. Vậy $y = v$. 

Kết hợp ba trường hợp, chúng ta có điều phải chứng minh.

\exercise Chứng minh các mệnh đề sau:
\begin{multicols}{2}
   \begin{enumerate}
      \item $\forall x, \forall y, \left(x = \{y\} \implies x \ni y\right)$;
      \item $\forall x, \forall y, \left(x = y \iff \{x\} = \{y\}\right)$;
      \item $\forall x, \forall y, \left(x \in y \iff \{x\} \subseteq y\right)$;
      \item $\forall x, \forall y, \left(\{x; y\} = \{x\} \iff x = y \right)$.
   \end{enumerate}
\end{multicols}

\solution

\setcounter{subexercise}{1}
\arabic{subexercise}. Giả sử $x = \{y\}$. Theo tính phản xạ của phép bằng nhau, $y = y$ với mọi $y$. Do vậy, $y \in \{y\}$. Khi đó, có $x = \{y\} \land y \in \{y\}$, sử dụng nguyên lí không thể phân biệt của các đối tượng đồng nhất, nên $y \in x$.

Do đó, với mọi $x$ và $y$, $x = \{y\} \implies x \ni y$. Điều phải chứng minh.

\stepcounter{subexercise}
\arabic{subexercise}. Trước hết, coi $x = y$ là đúng. Với mọi $z$, theo định nghĩa và tiên đề hợp thì
\begin{equation*}
   z \in \{x\} \iff z = x.
\end{equation*}
Tương tự, $z \in \{y\} \iff z = y$. Khi chúng ta kèm theo $x = y$, cần phải có $z = x \iff z = y$. Theo tính chất bắc cầu,
\begin{equation*}
   \forall z, z \in \{x\} \iff z \in \{y\}.
\end{equation*}
Vậy, $\forall x, \forall y, \left(x = y \implies \{x\} = \{y\}\right)$.

Ngược lại, coi $\{x\} = \{y\}$ đúng, dẫn đến $\forall w, \left(w \in \{x\} \iff w \in \{y\}\right)$. Ngoài ra, tương tự như trước, chúng ta cũng có khẳng định sau đúng với mọi $w$:
\begin{equation*}
   \begin{cases}
      w \in \{x\} \iff w = x \\
      w \in \{y\} \iff w = y
   \end{cases}.
\end{equation*}
Bắc cầu để được $\forall w, (w = x \iff w = y)$. Cụ thể hóa $w$ bởi $x$, chúng ta được $x = x \iff x = y$. Mà $x = x$ luôn đúng do tính phản xạ nên $x = y$. Vậy
\begin{equation*}
   \forall x, \forall y, \left(\{x\} = \{y\} \implies x = y\right).
\end{equation*}

Kết hợp hai chiều, điều phải chứng minh.

\stepcounter{subexercise}
\arabic{subexercise}. Thứ nhất, xét trường hợp đã có $\{x\} \subseteq y$. Khi này, theo định nghĩa tập con,
\begin{equation*}
   \forall w, \left(w \in \{x\} \implies w \in y\right).\\
\end{equation*}
Đổi $w$ bởi $x$, chúng ta có
\begin{equation*}
   x \in \{x\} \implies x \in y.
\end{equation*}
Mà $x \in \{x\}$ là hiển nhiên, nên $x \in y$.

Thứ hai, xét khi chúng ta có $x \in y$. Xét một phần tử $v$ thuộc $\{x\}$ bất kì. Dễ thấy $v = x$, và do đó $v \in y$. Vậy,
\begin{equation*}
   \forall v, v \in \{x\} \implies v \in y.
\end{equation*}
Dẫn đến, $\{x\} \subseteq y$.

Chúng ta đã chứng minh hai chiều --- $x \in y \impliedby \{x\} \subseteq y$ và $x \in y \implies \{x\} \subseteq y$ --- với mọi $x$ và $y$. Từ định nghĩa phép đẳng giá, chúng ta có điều phải chứng minh.